%-*-tex-*-
% Copyright Michael J. Ferguson, INRS-Telecommunications
% All rights reserved. 

%========= Autonumbering, Tagging and Autoreferencing =========

% These macros give the facility to symbolically refer to sections, equations,
% tables, and references and to automatically number them.
% Autonumber, automatically numbers sections, equations, ....
% Autoreferencing, causes \ref{tag} to drop the tag-value in the text.

% Although it is possible to have only autonumbering, referring to 
% sections or other parts is quite difficult without referencing 



% When autonumbering or autoreferencing is first turned on, 
% a tag file is brought in.
% Tags are written to a new file as they are generated.

\newif\ifautoreference
\newif\ifautonumber
\newif\ift@agfilein
\newwrite\t@agfile

\newtoks\inputtagfiles \inputtagfiles={\inputwithcheck {\jobname.tag} }
%\newtoks\inputtagfiles \inputtagfiles={\inputwithcheck {@\jobname.tag} }
% MSDOS version with renaming 
\def\g@ettagfile{\ift@agfilein
                           \else 
                             \the\inputtagfiles
                             \immediate\openout\t@agfile=\jobname.tag
                             \t@agfileintrue
                           \fi}
\def\autonumberingon{\g@ettagfile\autonumbertrue}
\def\autonumberingoff{\autonumberfalse}
\def\autoreferencingon{\g@ettagfile\autoreferencetrue}
\def\autoreferencingoff{\autoreferencefalse}

%-------- general tag setups --- new autonumbers
%#1 tagtype #2 pre
%tagrefformat will be called by \...tagrefformat and will be initially
% defined as \...tagrefformat --> \Prerefform\...numform\Postrefform
% and \...numform -->\the\...num
\def\Prerefform{} \def\Postrefform{}

%---- makes the tex commands for tag refformats -----
\def\m@aketagrefformat#1#2{\edef\nnxt{
    \noexpand\def\csname #1tagrefformat\endcsname{\noexpand\Prerefform #2\noexpand\Postrefform}}
       \nnxt}

% ----list of specials to be reset .... all are done together unless overidden
\newtoks\s@presetlist 



% ------ newautonum -------
% This is used to produce the commands necessary to have a tag that can be
% associated with a counter and to have the format styles for printing the
% numbers. They are used for such things a figures, tables, equations. 
% Three are actually defined. 

%  #1 is the name of then new autonum --- 3 or less characters
% since it might be the needed for list files
% if #1=xy, makes an \xynum, \xynumform \xytagrefformat
% sets \xynumform to \the\xynum
\def\newautonum#1{
                 \expandafter\n@ewcount\csname #1num\endcsname %allocates counter
                \edef\n@ext{{\global\expandafter\csname #1num\endcsname = 0}}
                 \expandafter\toks0\n@ext
                 \appendtoks{\s@presetlist}{\toks0}
                 \expandafter\let \csname #1tagrefformat\endcsname = \relax
                      %inhibits expansion later on
                 \edef\n@ext##1{\noexpand\def\csname auto#1num\endcsname
                     ##1{\noexpand\a@utotag{##1}{\csname #1num\endcsname
                         }{\csname #1tagrefformat\endcsname
                         }{\csname #1tagrefformat\endcsname}{##1}}}
                 \n@ext{##1} %makes \auto#1num
                 \m@aketagrefformat{#1}{\csname #1numform\endcsname}
                 \edef\n@ext{ %initializes \#1numform
                 \noexpand\def\csname #1numform\endcsname
                   {\noexpand\the\csname #1num\endcsname}}
                 \n@ext
                 }

%===== Allocates the autonums to eq, fig and tbl =======
\newautonum{eq}
\newautonum{fig}
\newautonum{tbl}
%--------------------
